% !TeX program = xelatex
\documentclass[11pt]{article}
\newcommand{\notesdir}{../../lecture-notes}
\providecommand{\notesdir}{.}

% -----------------------------------------------------------------------------
% Shared preamble for the CO5059 lecture notes and the quiz-review documents.
% Documents outside this folder define \notesdir (path to this folder) first.
% -----------------------------------------------------------------------------

% Page layout and typography
\usepackage[a4paper,portrait,margin=1in]{geometry}
\usepackage{fontspec}
\newfontfamily\arialfont{Arial}
\usepackage{microtype}
\usepackage{parskip}

% Mathematics and diagrams
\usepackage{amsmath,amssymb,amsthm,mathtools}
\usepackage{tikz-cd}
\usetikzlibrary{arrows.meta,positioning,shapes.geometric,shapes.arrows,calc,fit,backgrounds}

% Reusable styles for flowcharts and concept diagrams.
\tikzset{
  concept/.style={
    draw=brandprimary,
    fill=brandsecondary!8,
    rounded corners,
    minimum width=3.2cm,
    minimum height=0.9cm,
    align=center,
    font=\small,
  },
  diagram arrow/.style={-{Latex[length=2.5mm]},thick,brandprimary},
}

% Tables, figures, code, and lists
\usepackage{graphicx}
\graphicspath{{\notesdir/}}
\usepackage{enumitem}
\usepackage{booktabs}
\usepackage{array}
\usepackage{listings}

% Color and links
\usepackage[dvipsnames]{xcolor}
\definecolor{brandprimary}{HTML}{030391}
\definecolor{brandsecondary}{HTML}{1488D8}
\usepackage[
  colorlinks=true,
  linkcolor=brandprimary,
  citecolor=brandsecondary,
  urlcolor=brandsecondary
]{hyperref}

% Headers and footers
\usepackage{fancyhdr}
\setlength{\headheight}{36pt}
\pagestyle{fancy}
\fancyhf{}
\fancyhead[L]{%
  \parbox[c][0.42in][c]{0.42\textwidth}{%
    \raggedright\small\textbf{\coursecode}\\\coursename}}
\fancyhead[R]{%
  \parbox[c][0.42in][c]{0.42\textwidth}{%
    \raggedleft\includegraphics[height=0.42in]{assets/hcmut_logo.png}}}
\fancyfoot[L]{\small Semester \semester}
\fancyfoot[C]{\small\nouppercase{\leftmark}}
\fancyfoot[R]{\small Page \thepage}
\renewcommand{\sectionmark}[1]{\markboth{#1}{}}

% -----------------------------------------------------------------------------
% Course information -- change these values once per course
% -----------------------------------------------------------------------------
\newcommand{\universityname}{HCMC University of Technology}
\newcommand{\faculty}{Faculty of Computer Science and Engineering}
\newcommand{\coursecode}{CO5059}
\newcommand{\coursename}{Advanced System Architectures}
\newcommand{\semester}{261}
\newcommand{\instructor}{Tran Ngoc Thinh}
\newcommand{\instructoremail}{}% leave empty to hide
\newcommand{\notetaker}{Phan Nhat Minh}
\newcommand{\notetakeremail}{minh.phannhat95@hcmut.edu.vn}
\newcommand{\coursedescription}{Quantitative design and evaluation of computer
systems: performance and cost, instruction set architectures, pipelining and
its hazards, memory hierarchy, multiprocessors, superscalar/VLIW processors,
and AI \& security chips.}
\newcommand{\emaillink}[1]{%
  \href{mailto:#1}{\textcolor{brandprimary}{\texttt{#1}}}}

% Colored definitions, results, examples, proofs, and notes.
\input{\notesdir/colored-boxes}

% -----------------------------------------------------------------------------
% Course-specific notation
% -----------------------------------------------------------------------------
\newcommand{\CPI}{\mathrm{CPI}}

\lstdefinestyle{mips}{
  basicstyle=\ttfamily\small,
  keywordstyle=\bfseries\color{brandprimary},
  morekeywords={lw,sw,add,sub,and,or,slt,beq,bne,beqz,xor,nop,addi,addiu,addu,
    subu,li,la,jal,jr,jalr,sll,srl,sra,sllv,srlv,srav,andi,ori,xori,nor,slti,
    sltiu,sltu,lui,mult,multu,div,divu,mfhi,mflo,ll,sc,syscall,move,lh,lhu,lb,
    lbu,sh,sb,blez,bgtz,bltz,bgez,j,mul},
  showstringspaces=false,
  columns=fullflexible,
  keepspaces=true,
}
\lstset{style=mips}

% Creates a consistently formatted lecture heading.
% Usage: \lecture{number}{title}{date}
\newcommand{\lecture}[3]{%
  \section{Lecture #1: #2}
  \label{lec:#1}
  \begin{center}
    \small\textit{#3}
  \end{center}
}


% Cover page + contents page.  Usage: \makecover{Big title}{Description}
\newcommand{\makecover}[2]{%
\pagenumbering{roman}
\begin{titlepage}
  \thispagestyle{empty}

  \begin{center}
    \includegraphics[width=0.82\textwidth]{assets/hcmut_banner_en.png}

    \vspace{0.35cm}
    {\arialfont\large\bfseries\color{brandprimary}\faculty\par}
  \end{center}

  \vfill
  \begin{center}
    {\Huge\bfseries #1\par}
    \vspace{0.55cm}
    {\LARGE\bfseries\color{brandprimary}\coursename\par}
    \vspace{0.2cm}
    {\arialfont\large\bfseries\color{brandsecondary}\coursecode\par}

    \vspace{0.65cm}
    {\color{brandsecondary}\rule{0.38\textwidth}{1.2pt}}
    \vspace{0.65cm}

    {\large\textbf{Instructor:} \instructor
      \ifx\instructoremail\empty\else\ {\small(\emaillink{\instructoremail})}\fi\par}
    \vspace{0.3cm}
    {\large\textbf{Notes by:} \notetaker\
      {\small(\emaillink{\notetakeremail})}\par}
    \vspace{0.3cm}
    {\large\textbf{Semester:} \semester\par}

    \vspace{0.8cm}
    \begin{minipage}{0.72\textwidth}
      \centering\itshape\color{black!70}
      #2
    \end{minipage}
  \end{center}

  \vfill
  \begin{center}
    \small\color{black!60}Last updated: \today
  \end{center}
\end{titlepage}

% -----------------------------------------------------------------------------
% Table of contents -- kept on its own page
% -----------------------------------------------------------------------------
\clearpage
\setcounter{page}{2}
\thispagestyle{empty}
\hypersetup{linkcolor=black}
\setcounter{tocdepth}{2}
\tableofcontents
\hypersetup{linkcolor=brandprimary}
\clearpage
\pagenumbering{arabic}
}


\newcommand{\ex}[3]{\subsection*{Exercise #1 \normalfont\small(book p.~#2, PDF p.~#3)}%
  \addcontentsline{toc}{subsection}{Exercise #1}}

\begin{document}
\sloppy
\makecover{Textbook Exercises: Pipelining}{Hennessy \& Patterson,
\emph{Computer Architecture: A Quantitative Approach}, 5th ed., Appendix~C
(Pipelining: Basic and Intermediate Concepts) exercises C.1--C.14. A map of
every exercise with its page and relevance to Lecture~3, worked solutions for
the ones our lectures cover, and hints for the rest.}

% ============================================================================
\section{How to Use This File}

\begin{itemize}[leftmargin=*]
  \item The exercise text is \textbf{not} copied here -- open the book
  (\texttt{references/Computer Architecture - A Quantitative Approach.pdf})
  at the page given. ``Book p.'' is the printed page (Appendix~C pages are
  numbered C-1, C-2, \dots); ``PDF p.'' is the page in the PDF viewer.
  \item The book's pipeline is the same 5-stage MIPS pipeline as our lecture
  (it writes 64-bit opcodes such as \texttt{LD}, \texttt{DADDI},
  \texttt{BNEZ}).
  \item Unless an exercise says otherwise, the solutions assume what the book
  does in Section~C.2: registers are written in the first half and read in
  the second half of a cycle, and branches are resolved in \textbf{ID}, so
  ``flushing'' costs one cycle (Fig.~C.11). In diagrams, \texttt{stl} means
  a stall cycle.
\end{itemize}

% ============================================================================
\section{Exercise Map}

\begin{center}
\small
\renewcommand{\arraystretch}{1.15}
\begin{tabular}{@{}lll>{\raggedright\arraybackslash}p{5.4cm}l@{}}
\toprule
Ex. & Book p. & PDF p. & Topic & Priority \\
\midrule
C.1 & C-82/83 & 729--730 & Loop: dependences, timing without/with forwarding, branch schemes, 10-stage pipeline, CPI & \textbf{High} \\
C.2 & C-83 & 730 & Branch penalties in a 4-deep and a 15-deep pipeline & \textbf{High} \\
C.3 & C-83 & 730 & Unbalanced stages: clock, CPI with stalls, speedup & \textbf{High} \\
C.4 & C-84 & 731 & Branch resolved in ID creates a data hazard & \textbf{High} \\
C.5 & C-84/85 & 731--732 & Pipeline for a register-memory ISA: adders, ports, forwarding, hazards & Medium \\
C.6 & C-85/86 & 732--733 & Adding register-memory ALU ops to the 5-stage pipeline & Medium \\
C.7 & C-86 & 733 & Deeper pipeline: faster clock vs.\ more stalls & \textbf{High} \\
C.8--C.11 & C-86 & 733 & R4000 forwarding/hazard tables; FP pipeline tables & Skip (later topic) \\
C.12--C.13 & C-86/88 & 733--735 & DAXPY on FP pipeline; scoreboarding & Skip (later topic) \\
C.14 & C-88 & 735 & Multicycle FP units: structural and WAW hazards & Skip (later topic) \\
\bottomrule
\end{tabular}
\end{center}
``Later topic'': multicycle FP pipelines, the R4000 and scoreboarding
(Sections~C.5--C.7) go beyond our Lecture~3; they belong with the
instruction-level-parallelism lectures.

% ============================================================================
\section{Worked Solutions: High Priority}

\ex{C.1}{C-82/83}{729--730}
\emph{Asks:} for the loop \texttt{LD R1; DADDI R1; SD R1; DADDI R2; DSUB R4;
BNEZ R4} (99 iterations, since R3 $=$ R2 $+$ 396 and R2 grows by 4): list
the dependences; time the loop without forwarding (flush on branch), with
forwarding (predict not taken / predict taken), and on a 10-stage pipeline;
then clock cycles and CPI.
\begin{solution}
\textbf{(a)} RAW dependences: R1 from \texttt{LD} to \texttt{DADDI R1};
R1 from \texttt{DADDI R1} to \texttt{SD}; R2 from \texttt{DADDI R2} to
\texttt{DSUB}; R4 from \texttt{DSUB} to \texttt{BNEZ}; and the loop-carried
R2 from \texttt{DADDI R2} to the next iteration's \texttt{LD} and
\texttt{SD}. (Also WAR: \texttt{LD}/\texttt{SD} read R2 before
\texttt{DADDI R2} writes it.)

\textbf{(b)} No forwarding; a consumer's ID can be no earlier than the
producer's WB:
\begin{lstlisting}[basicstyle=\ttfamily\scriptsize,keywords={}]
                 c1  c2  c3  c4  c5  c6  c7  c8  c9  c10 c11 c12 c13 c14 c15 c16 c17 c18
LD    R1,0(R2)   IF  ID  EX  MEM WB
DADDI R1,R1,#1       IF  ID  stl stl EX  MEM WB
SD    R1,0(R2)           IF  stl stl ID  stl stl EX  MEM WB
DADDI R2,R2,#4                       IF  stl stl ID  EX  MEM WB
DSUB  R4,R3,R2                                   IF  ID  stl stl EX  MEM WB
BNEZ  R4,Loop                                        IF  stl stl ID  stl stl EX  MEM WB
LD (next iter.)                                                              IF
\end{lstlisting}
The branch is resolved in ID in cycle 15, so the next iteration's
\texttt{LD} is fetched in cycle 16: 15 cycles per iteration. Total $=
98\times15 + 18 = 1488$ cycles.

\textbf{(c)} Full forwarding, predict not taken. Only the load-use pair
stalls, plus one stall because \texttt{BNEZ} needs R4 in ID but
\texttt{DSUB} produces it at the end of EX (exactly the problem of C.4):
\begin{lstlisting}[basicstyle=\ttfamily\scriptsize,keywords={}]
                 c1   c2   c3   c4   c5   c6   c7   c8   c9   c10  c11  c12  c13  c14
LD    R1,0(R2)   IF   ID   EX   MEM  WB
DADDI R1,R1,#1        IF   ID   stl  EX   MEM  WB
SD    R1,0(R2)             IF   stl  ID   EX   MEM  WB
DADDI R2,R2,#4                       IF   ID   EX   MEM  WB
DSUB  R4,R3,R2                            IF   ID   EX   MEM  WB
BNEZ  R4,Loop                                  IF   ID   stl  EX   MEM  WB
LD (next iter.)                                               IF   ID   EX   MEM  WB
\end{lstlisting}
\texttt{BNEZ} resolves in cycle 9; the predicted fall-through instruction is
squashed and the target \texttt{LD} is fetched in cycle 10: 9 cycles per
iteration. Total $= 98\times9 + 12 = 894$ cycles.

\textbf{(d)} Predict taken gives the same 894 cycles: in this pipeline the
target address is also computed in ID, so a taken branch cannot be
fetched any earlier (as noted in our lecture).

\textbf{(e)} 10 stages (IF1 IF2 ID1 ID2 EX1 EX2 M1 M2 WB1 WB2), forwarding
from the end of a stage pair to the start of the pair that needs it. If
instruction $p$ starts at cycle $s_p$, a dependent ALU instruction must start
at $s_p+2$ or later, and a dependent use of a load at $s_p+4$ or later.
Assuming the branch target is known after ID2 and the condition is tested in
EX1: \texttt{LD} 0, \texttt{DADDI} 4, \texttt{SD} 5, \texttt{DADDI R2} 6,
\texttt{DSUB} 8, \texttt{BNEZ} 10, next \texttt{LD} at 14. That is 14 cycles
per iteration, about $98\times14 + 20 = 1392$ cycles in total. (Different
assumptions about where the branch resolves change this number; state
yours.)

\textbf{(f)} 5-stage: $0.8 + 0.1 = 0.9$~ns. 10-stage: $0.4 + 0.1 =
0.5$~ns.

\textbf{(g)} 5-stage: 9 cycles per 6 instructions, $\CPI = 1.5$, average
time per instruction $=1.5\times0.9 = 1.35$~ns. 10-stage: $14/6 = 2.33$
cycles, $2.33\times0.5 = 1.17$~ns. The deeper pipeline wins only modestly:
its clock is 1.8$\times$ faster but its CPI is 1.56$\times$ worse.
\end{solution}

\ex{C.2}{C-83}{730}
\emph{Asks:} with 15\% conditional branches (60\% taken) and 1\% jumps, how
much faster the machine would be without branch hazards, for a 4-deep and
a 15-deep pipeline.
\begin{solution}
If a branch is resolved at the end of stage $r$ and only stage 1 can proceed
regardless, the penalty is $r-1$ cycles.
\begin{enumerate}[label=(\alph*),leftmargin=*]
  \item Jumps: $2-1 = 1$; conditional: $3-1 = 2$.
  $\CPI = 1 + 0.01\times1 + 0.15\times2 = 1.31$, so $1.31\times$ faster
  without hazards.
  \item Jumps: 4; conditional: 9. $\CPI = 1 + 0.01\times4 + 0.15\times9 =
  2.39$, so $2.39\times$ faster.
\end{enumerate}
Deep pipelines make branches dramatically more expensive -- the motivation
for good branch prediction. (If you assume only taken branches pay, the
numbers become $1 + 0.01 + 0.15\times0.6\times2 = 1.19$ and $1 + 0.04 +
0.15\times0.6\times9 = 1.85$; say which you assume.)
\end{solution}

\ex{C.3}{C-83}{730}
\emph{Asks:} a 7~ns single-cycle machine split into stages of 1, 1.5, 1,
2, 1.5~ns with 0.1~ns register delay: pipelined clock, CPI with a stall
every 4 instructions, speedup, and speedup with infinitely many stages.
\begin{solution}
\leavevmode
\begin{enumerate}[label=(\alph*),leftmargin=*]
  \item $2 + 0.1 = 2.1$~ns (slowest stage plus register).
  \item $\CPI = 1 + 1/4 = 1.25$.
  \item $7/(1.25\times2.1) = 7/2.625 = 2.67$.
  \item With infinitely many stages the logic per stage goes to 0 and the
  clock approaches the 0.1~ns register delay: speedup $\to 7/0.1 = 70$
  (or $7/(1.25\times0.1) = 56$ if the 1-in-4 stalls remain). Register
  overhead is the limit on how deep a pipeline is worth making.
\end{enumerate}
\end{solution}

\ex{C.4}{C-84}{731}
\emph{Asks:} a code snippet where resolving the branch in ID causes a data
hazard even with forwarding.
\begin{solution}
\begin{lstlisting}
DADD  R1, R2, R3
BEQZ  R1, Target
\end{lstlisting}
\texttt{BEQZ} needs R1 in its ID stage (cycle 3), but \texttt{DADD} produces
it only at the end of its EX (cycle 3): one stall even with forwarding. A
load is worse:
\texttt{LD R1,0(R2); BEQZ R1,Target} needs 2 stalls, because the data exists
only at the end of MEM. Moving the branch earlier saves branch cycles but
creates these new data stalls.
\end{solution}

\ex{C.7}{C-86}{733}
\emph{Asks:} 5-stage pipeline (1~ns, a data stall every 5 instructions) vs.\
12-stage (0.6~ns, 3 stalls every 8 instructions); 20\% branches with 5\%
mispredicted, penalties 2 and 5 cycles.
\begin{solution}
\leavevmode
\begin{enumerate}[label=(\alph*),leftmargin=*]
  \item Data hazards only: $\CPI_5 = 1 + 1/5 = 1.2$, $\CPI_{12} = 1 + 3/8 =
  1.375$. Time per instruction: $1.2\times1 = 1.2$~ns vs.\
  $1.375\times0.6 = 0.825$~ns. Speedup $=1.2/0.825 = 1.45$.
  \item Mispredictions per instruction $=0.2\times0.05 = 0.01$:
  $\CPI_5 = 1.2 + 0.01\times2 = 1.22$, $\CPI_{12} = 1.375 + 0.01\times5 =
  1.425$. (Speedup is now $1.22/(1.425\times0.6) = 1.43$.)
\end{enumerate}
\end{solution}

% ============================================================================
\section{Hints: Medium Priority}

\ex{C.5}{C-84/85}{731--732}
\emph{Asks:} a 6-stage pipeline (IF RF ALU1 MEM ALU2 WB) for a
register-memory ISA: count adders and register/memory ports, design the
forwarding, list remaining hazards and control hazards.
\begin{solution}
Walk an instruction mix through the stages cycle by cycle and look for the
worst case in a single cycle:
\begin{itemize}[leftmargin=*]
  \item Adders: PC increment (IF), address/branch-target add (ALU1),
  operation (ALU2) -- several can be busy at once.
  \item Ports: RF reads up to 2 (or 3 for a store with a register-memory
  form), WB writes 1, memory: 1 instruction read and 1 data access per cycle.
  \item Forwarding: results appear at the end of ALU2 or MEM and are needed
  at the start of ALU1 (addresses) or ALU2 (operands); consider forwarding
  across one intervening instruction too.
  \item A result needed for an address in ALU1 one instruction later cannot
  be forwarded in time: those become stalls (address-generation
  interlocks). Branches compare in ALU2, so the branch penalty is several
  cycles.
\end{itemize}
\end{solution}

\ex{C.6}{C-85/86}{732--733}
\emph{Asks:} reorder the 5 stages so ALU operations can take a memory operand
addressed by a register only; new forwarding paths, new hazards, and effects
on instruction count and CPI.
\begin{solution}
Put MEM before EX (IF ID MEM EX WB): with register-indirect addressing no
address calculation is needed, so memory can be read before the ALU. New
paths: forward results to the MEM stage (addresses) as well as to EX. New
hazard: an ALU result used as an address by the next instruction must wait
(the address is needed in MEM, which now comes \emph{before} EX).
Instruction count drops when a load$+$ALU pair becomes one instruction, but
rises when displacement addressing must be replaced by an explicit add. CPI
changes through these new address hazards and the removed load-use stalls.
\end{solution}

\end{document}
